\section{Experimental Setup and Additional Results}
\label{app:setup}

\subsection{Models and serving}
All agents and generators run locally via vLLM 0.31.0 (OpenAI-compatible
API, \texttt{--enforce-eager}, temperature 0 for agents and the judge,
temperature 0.7 with a fixed seed for the generator) on RTX 4090 GPUs
(24\,GB; two cards for the original battery, one for the revision
rounds): Qwen2.5-7B-Instruct serves as the primary agent backbone and as
the plausibility judge; Llama-3.1-8B-Instruct serves as the contamination
generator and as the secondary agent backbone. Qwen2.5-3B/1.5B are used in
the scale ablation. Every LLM call is logged with its token counts; agents
cap episodes at ten steps, forced-choice probes at six.

\subsection{Official artifacts}
We use WebShop's official \texttt{items\_shuffle.json} (1{,}181{,}430
products after key cleaning and deduplication), \texttt{items\_ins\_v2.json},
and \texttt{items\_human\_ins.json} (10{,}136 human instructions), with the
benchmark's own loader, goal builder, and reward function applied
unmodified; the only semantic addition is best-variant scoring of the
chosen listing (\S\ref{sec:method-env}).

\subsection{Contamination pool}
The pool covers 100 tasks (the official goal list under seed 2027), three
families, four listings per family per task (up to 480 per task set), with a
plausibility floor at judge score~4 (regenerate once, then drop). Synthetic
economics are assigned programmatically: price = gold price $\times$
$\mathrm{U}(0.75,0.95)$ (deterministic per listing), rating
$\mathrm{U}(4.6,4.9)$, review counts in $[400,3000)$, and reviews in packs
of 20 (reviews family) or 6 (others). An archived pre-fix pool with a
\$50 price anchor is released as an ablation.

\subsection{Gate outcomes}
\label{app:gate}
P0: success $0.65$ at $\varepsilon=0$, zero malformed actions, 18/20 tasks
exactly reproducible under greedy decoding (two flips attributable to
continuous-batching numerics). P1 (final, corrected pool): $\sigma=0.377$
[0.324, 0.433] over 300 duels; per family: clone $0.523$, geo $0.351$,
reviews $0.287$. We report this gate honestly: as pre-registered, P1 asked
for an overall susceptibility of $\sigma\ge0.60$, which the final pool did
\emph{not} meet --- agents on average \emph{reject} machine-optimized
pages. The gate's intent, a decision-level susceptibility signal, was met
in modified form by the strong family differentiation (clones at chance
versus gloss significantly rejected), which redirected the study from
``uniform gullibility'' toward channel-specific economic lures; every
measurement after the probe was pre-registered against this revised
premise, and the original gate text is preserved in the repository. P2 (gate run, 40 paired tasks): success $0.525\to0.450$,
regret $0.338\to0.398$, confidence $0.783\to0.862$ at $\varepsilon=0.75$
versus $0$. The main dose--response curve ($N{=}100$ per dose)
yields: success
$0.590\to0.360$ (McNemar $p=1.3\times10^{-4}$), regret $0.262\to0.443$
(paired $\Delta=+0.181$, CI $[+0.095,+0.264]$), confidence $0.822\to0.875$,
ECE $0.249\to0.496$; P4E replicates the endpoints under a second
contamination seed (pooled McNemar $p=1.2\times10^{-8}$). Two earlier
P1/P2 runs were invalidated by a rendering defect (machine-readable review
dictionaries leaking into rendered pages) and a price-anchoring bug; both
are documented with their data in the repository.

\subsection{Power and multiple comparisons}
\label{app:power}
The paired regret difference at $\varepsilon=0.75$ has task-level SD
$\approx0.43$; at $N=100$ per dose the endpoint McNemar test detects the
observed success drop ($p=1.3\times10^{-4}$), and pooling the second seed
tightens it further ($p=1.2\times10^{-8}$); the paired regret difference is
significant with a bootstrap CI excluding zero. Intermediate doses are
reported as effect sizes with cluster-bootstrap CIs rather than pairwise
p-values. We do not over-claim single small deltas; the interpretation
rests on the joint pattern across outcomes (success, regret, purchase
share, confidence, calibration) and its replication across families,
backbones, scales, and defenses.

\paragraph{Family-wise correction.} Applying Holm--Bonferroni across the
eleven confirmatory tests reported in the paper leaves the two headline
dose contrasts unambiguous (pooled success: adjusted $p=1.3\times10^{-7}$;
full-parity probe: adjusted $p=1.0\times10^{-3}$) and the $10\%$-discount
threshold (adjusted $p=0.023$); the remaining channel comparisons
(price parity, rating parity, provenance-vs-none, stealth recovery, the
$s{=}25\%$ envelope point) survive uncorrected but not Holm-corrected
(adjusted $p\in[0.072,0.078]$). We therefore read the decomposition as a convergent joint pattern ---
two channels of nearly equal magnitude, an additive account, and an
independent threshold law reaching the same conclusion --- not as a
collection of individually significant contrasts.

\paragraph{Equivalence at saturation.} The pricing-threshold claim that
``deeper discounts add nothing'' is supported by two-proportion tests
against the $10\%$ point ($20\%$: $p=0.45$; $30\%$: $p=0.76$; $20\%$ vs.\
$30\%$: $p=0.65$; $n{=}88$ per arm). These are non-significance results with modest power; we phrase the
claim as saturation of the observed response, not as a strict
equivalence bound.

\paragraph{Structural detector vs.\ oracle.} With $n{=}40$ paired tasks,
the minimal detectable regret difference at two-sided $\alpha=0.05$ and
$80\%$ power is roughly $0.11$ --- larger than the observed
oracle--struct gap ($0.171$ vs.\ $0.250$). We conclude only that structural
filtering recovers at least the full regret excess ($0.398\to0.250$,
below the uncontaminated $0.338$); we do not claim equality with the
oracle.

\subsection{Real-world attack-surface mining}
\label{sec:app-real}
Section~\ref{sec:res-realworld} grounds the tested discount regime in
real catalogs. This appendix documents the three sources and a text
provenance check.

\paragraph{Sources and pairing rules.}
\emph{ProMap}~\cite{mackov2023promap} provides human-verified
Walmart--Amazon pairs split into true matches, close non-matches
(different products sharing brand and product line with
near-identical titles), and medium non-matches. We compute the
cheaper-side discount per pair with both prices listed
(330/509 close non-matches, 337/509 matches, 364/537 medium).
\emph{WebMall}~\cite{peeters2025webmall} offers are grouped across its
four shops by exact normalized name (244 groups, 505 offers; the
GTIN field is empty in the distribution). \emph{WebShop}~\cite{yao2022webshop}
catalog items (80{,}000 sampled from the 1.18M-item real scrape) are
blocked by first brand token and category; within blocks, pairs with
title-token Jaccard in $[0.80,0.95)$ and distinct ASINs are lookalikes,
Jaccard $\geq0.95$ near-identical (232 pairs). Discount statistics
appear in Section~\ref{sec:res-realworld}; per-source values are
 archived in the repository (\texttt{realdata/real\_stats.json}).
Close non-match pairs are the real-world analogue of our clone
channel; identical-name groups across shops play the near-identical
control role, where the cheaper offer is the correct purchase.

\paragraph{Can provenance be detected from text alone?}
A natural remedy would be a text-side detector that flags synthetic
listings at browsing time. We calibrate the standard GPT-2
perplexity signal~\cite{radford2019gpt2} on 250 human-written
WebShop descriptions and score our synthetic pools: separation is
strong (AUC $0.912$ against a Mann--Whitney null; synthetic median
perplexity $\approx22$ vs.\ $\approx55$ human), i.e., our generated
listings are identifiable \emph{in isolation} by a commodity
detector. Two conclusions follow. First, susceptibility is a
\emph{decision-level} failure: the agent buys listings it could in
principle distinguish, which is why structural provenance guarantees
rather than better perception are the operative defense
(Section~\ref{sec:res-defense}). Second, we deliberately do
\emph{not} estimate the share of AI-generated text in current real
catalogs from this signal: applying the same threshold to WebMall
offer text flags a large fraction, but the estimate is dominated by
genre and length confounds (flagging falls from 93.6\% to 55.3\%
under length matching alone), and no ground truth exists for
real-offer provenance. Recent product-matching
benchmarks~\cite{steiner2026billigerde} likewise treat provenance as
unknown. Text-only prevalence claims are left to work with
calibrated multi-signal detectors and verified labels.
